\begingroup
\setlength{\parindent}{0pt}
\setlength{\parskip}{0.8em}

\section*{Sammendrag}
Rapporten omhandler forbedring av det ugraderte nettverket
til bataljonsstridsgruppen i Rukla, Litauen. Utgangspunktet
er et nettverk med mangelfull dokumentasjon, begrenset
kontroll med tilkoblet utstyr og manglende skille mellom
brukergrupper og tjenester. Dette medfører risiko for
uønsket tilgang og driftsavbrudd som kan ramme telefoni,
PA-anlegg, kontorarbeid og velferdstjenester.

Formålet er å utarbeide et overordnet nettverksdesign som
reduserer sikkerhetsrisikoen og ivaretar eksisterende
funksjonalitet. Løsningen skal ta utgangspunkt i tilgjengelig
utstyr, uten større nyanskaffelser eller endringer i
forbindelsene levert av det litauiske forsvaret. Drift og
bruk skal heller ikke kreve mer ressurser enn i dag.

Det faglige grunnlaget omfatter logisk oppdeling av
nettverket, kontroll av trafikk mellom sikkerhetssoner og
beskyttelse av kommunikasjon mellom bygg. Videre behandles
tiltak mot uautoriserte nettverkstjenester, falske
nettverksadresser og feilkoblinger. Sikker fjernadministrasjon,
sentral logginnsamling, felles tidsreferanse og overvåkning
av nettverkstrafikk inngår som grunnlag for bedre kontroll,
feilsøking og oppfølging av sikkerhetshendelser.

Aktuelle designvalg skal vurderes gjennom laboratorietester
og drøfting av tekniske og praktiske hensyn. Arbeidet skal
gi S6 et beslutningsgrunnlag for eventuell rekonfigurasjon
av produksjonsnettet. Dokumentasjonen skal samtidig støtte
overføring av kunnskap til fremtidige kontingenter og bidra
til at nettverket kan driftes og videreutvikles på en
forsvarlig måte.

\par
\endgroup